% =====================================================================
%  preamble/style.tex — margenes, encabezados y apariencia
% =====================================================================

% --- Geometria de pagina ---------------------------------------------
\usepackage[
  a4paper,
  top=2.5cm, bottom=2.5cm,
  left=3cm,  right=2.5cm,
  headheight=14pt,
]{geometry}

% --- Espaciado -------------------------------------------------------
\usepackage{setspace}
\onehalfspacing
\setlength{\parskip}{0.5em}

% --- Encabezados y pies ----------------------------------------------
\usepackage{fancyhdr}
\pagestyle{fancy}
\fancyhf{}
\fancyhead[L]{\small\nouppercase{\leftmark}}
\fancyhead[R]{\small\thepage}
\renewcommand{\headrulewidth}{0.4pt}

% --- Paleta de color -------------------------------------------------
\definecolor{acento}{HTML}{1F4E79}
\definecolor{grisclaro}{HTML}{F5F5F5}
\definecolor{comentario}{HTML}{6A737D}

% --- Titulos de seccion ----------------------------------------------
\usepackage[explicit]{titlesec}
\titleformat{\section}
  {\Large\bfseries\color{acento}}{\thesection}{1em}{#1}
\titleformat{\subsection}
  {\large\bfseries}{\thesubsection}{1em}{#1}
\titlespacing*{\section}{0pt}{2.5ex plus 1ex minus .2ex}{1.5ex plus .2ex}

% --- Estilo de los bloques de codigo ---------------------------------
\lstdefinestyle{base}{
  backgroundcolor=\color{grisclaro},
  commentstyle=\color{comentario}\itshape,
  keywordstyle=\color{acento}\bfseries,
  stringstyle=\color{teal},
  basicstyle=\ttfamily\footnotesize,
  numberstyle=\tiny\color{comentario},
  numbers=left,
  numbersep=8pt,
  breaklines=true,
  breakatwhitespace=false,
  showstringspaces=false,
  tabsize=2,
  frame=single,
  rulecolor=\color{comentario!40},
  captionpos=b,
  extendedchars=true,
  literate=%
    {á}{{\'a}}1 {é}{{\'e}}1 {í}{{\'i}}1 {ó}{{\'o}}1 {ú}{{\'u}}1
    {Á}{{\'A}}1 {É}{{\'E}}1 {Í}{{\'I}}1 {Ó}{{\'O}}1 {Ú}{{\'U}}1
    {ñ}{{\~n}}1 {Ñ}{{\~N}}1 {ü}{{\"u}}1 {¿}{{?`}}1 {¡}{{!`}}1
}
\lstset{style=base}

% --- Entornos de teorema --------------------------------------------
\theoremstyle{definition}
\newtheorem{definicion}{Definicion}[section]
\newtheorem{ejemplo}{Ejemplo}[section]

\theoremstyle{plain}
\newtheorem{teorema}{Teorema}[section]
\newtheorem{lema}[teorema]{Lema}

\theoremstyle{remark}
\newtheorem*{nota}{Nota}
